\section{Introduction}

Plastic waste generation has become one of the most critical environmental challenges of the twenty-first century, with global waste volumes expected to escalate substantially in the coming decades. Recent studies estimate that global plastic production exceeded 400~Mt in 2022~\cite{houssini2025complexities} and will rise to approximately 735.4--832.4~Mt by 2050~\cite{dokl2024global}. Correspondingly, the OECD projects that annual plastic waste generation will increase from 353~Mt in 2019 to 1{,}014~Mt by 2060~\cite{oecd2022PlasticWaste2060}. The mismanagement of plastic waste further intensifies environmental degradation: between 60 and 99~Mt of mismanaged plastic waste were inadequately disposed into the environment in 2015, a figure projected to nearly triple to 155–265~Mt per year by 2060~\cite{lebreton2019future}. These trends underscore the urgent need for robust, data-driven forecasting tools to support proactive planning and evidence-based intervention strategies.

In response to these situations, recent research highlights the pivotal role of policy interventions in mitigating future plastic waste mismanagement. Pottinger et al.~\cite{pottinger2024pathways} demonstrate, through a comprehensive global systems model, that a wide range of policy levers—including minimum recycled content mandates, enhanced recycling collection rates, packaging reuse requirements, consumption taxes, and targeted investments in waste and recycling infrastructure—can substantially reduce mismanaged plastic waste by 2050. The effectiveness of the interventions requires data-driven insights into the temporal dynamics of plastic waste generation. However, current decision-making frameworks used by governments, NGOs, and industry largely rely on static estimates and incomplete datasets, limiting proactive planning and the identification of high-impact interventions. The Alliance to End Plastic Waste (AEPW) has invested in the Plastics Recovery Insights Steering Model (PRISM) data ecosystem to address these gaps \cite{edie2023IBMPrismPlatform}. This eceosystem leverages large language models (LLMs) to extract various data points from unstructured sources, including government reports, academic publications, and industry datasets, to create a comprehensive repository of plastic waste-related information.

Building on PRISM, this paper presents AUDA as an analytical platform for environmental data modeling. AUDA migrates validated PRISM records into a normalized country-year database and provides reusable workflows for constructing modeling datasets, filtering anomalous observations, reconstructing missing values, forecasting temporal trends, and evaluating model performance. Using AUDA, we conduct a set of experiments on plastic waste and related driver variables, including demographic, socioeconomic, waste-management, resin-consumption, and polymer-specific indicators.

% Codex
% TODO: rewrite this
This paper also introduces a variant of the one-standard-error rule for hyperparameter selection. The standard one-standard-error rule favors simpler models whose validation error is statistically close to the best observed validation error. In sparse environmental datasets, however, the tolerance used by this rule can strongly affect which configuration is selected. We therefore introduce a selection-efficiency coefficient, \(c_{\mathrm{SE}}\), that controls how much validation-error uncertainty is allowed when selecting simpler configurations. This coefficient makes the trade-off between predictive performance and model simplicity explicit, and allows the selection rule to be tested for sensitivity and robustness.

% Codex
% TODO: rewrite this
The paper makes two main contributions. First, it provides an empirical benchmark for analyzing sparse environmental data, covering correlation analysis, feature importance, reconstruction, temporal forecasting, and multivariate forecasting across datasets ranging from country-specific time series to pooled cross-country driver records. Second, it evaluates the proposed selection-efficiency variant of the one-standard-error rule and studies how robust the selected hyperparameters are under different tolerance and complexity-ordering assumptions. Together, these contributions connect the practical problem of plastic waste data analysis with a reproducible model-selection procedure designed for small, noisy, and heterogeneous environmental datasets.
